\section{共享年表之一：东迁、齐霸与晋楚相遇}

\begin{event}{公元前770年}{平王东迁洛邑}{可靠纪年}{洛邑}\eventanchor{SA-0770-EAST}{春秋·东迁}
\term{这一年发生了什么}：平王迁都洛邑，东周开始。
\term{后来改变了什么}：周王室从关中资源中心退到中原礼仪中心。此后任何诸侯都可以“尊王”，却很少再接受王室直接支配。
\end{event}

\begin{event}{公元前707年}{繻葛之战}{可靠纪年}{繻葛}\eventanchor{SA-0707-XUGE}{春秋·繻葛之战}
\term{这一年发生了什么}：郑庄公与周桓王交战，周军失利，桓王受伤。
\term{后来改变了什么}：这场战役并未废除周天子的名义，却让诸侯和后人更清楚地看到：王室的军事权威已经不能自动兑现。
\end{event}

\begin{event}{公元前685年}{齐桓公即位，管仲入齐}{可靠纪年}{齐}\eventanchor{SA-0685-QIHUAN}{春秋·齐桓公即位}
\term{这一年发生了什么}：齐国在内乱后重建权力，齐桓公任用管仲。
\term{事情为什么重要}：齐能组织会盟，靠的不只是个人声望，还靠滨海资源、内政整顿和把诸侯利益装进“尊王攘夷”语言的能力。
\end{event}

\begin{event}{公元前632年}{城濮之战}{可靠纪年}{城濮}\eventanchor{SA-0632-CHENGPU}{春秋·城濮之战}
\term{这一年发生了什么}：晋、楚在城濮交战，晋文公取胜。
\term{后来改变了什么}：齐的早期霸权之后，晋成为中原联盟的关键组织者；楚则显示南方大国已无法被排除在“天下”之外。
\end{event}
